%% ThreadLearn: Research Paper
%% Template: article onecolumn (clean layout for review)
%% Target venue: ASPLOS 2027 or ICLR 2026

\documentclass[11pt]{article}

%% ── Geometry ────────────────────────────────────────────────────────────────
\usepackage[a4paper, top=1in, bottom=1in, left=1.25in, right=1.25in]{geometry}

%% ── Packages ────────────────────────────────────────────────────────────────
\usepackage{booktabs}
\usepackage{multirow}
\usepackage{listings}
\usepackage{xcolor}
\usepackage{graphicx}
\usepackage{microtype}
\usepackage{enumitem}
\usepackage{fancyhdr}
\usepackage{amssymb}
\usepackage{amsmath}
\usepackage[hidelinks]{hyperref}
\usepackage{setspace}
\usepackage{caption}

%% ── Typography ───────────────────────────────────────────────────────────────
\setlength{\parskip}{4pt}
\setlength{\parindent}{0pt}
\captionsetup{font=small, labelfont=bf}

%% ── Header / footer ─────────────────────────────────────────────────────────
\pagestyle{fancy}
\fancyhf{}
\fancyhead[L]{\small\textit{ThreadLearn --- Under Review}}
\fancyhead[R]{\small\thepage}
\renewcommand{\headrulewidth}{0.4pt}

%% ── Code listing style ──────────────────────────────────────────────────────
\lstdefinestyle{codebase}{
  basicstyle=\small\ttfamily,
  keywordstyle=\color{blue!70!black}\bfseries,
  commentstyle=\color{gray!70},
  stringstyle=\color{green!50!black},
  numberstyle=\tiny\color{gray},
  numbers=left,
  numbersep=8pt,
  frame=single,
  framesep=4pt,
  breaklines=true,
  showstringspaces=false,
  tabsize=2,
  xleftmargin=12pt,
}
\lstdefinestyle{js}{style=codebase, language=JavaScript}
\lstdefinestyle{py}{style=codebase, language=Python}

%% ── Title ───────────────────────────────────────────────────────────────────
\title{%
  \large\textbf{ThreadLearn: Combining Static Race Detection with}\\[2pt]
  \large\textbf{Hindsight Chain-of-Thought Fine-Tuning}\\[2pt]
  \large\textbf{for Concurrency Bug Remediation}\\[8pt]
  \normalsize\textit{Under Review --- Target Venue: ASPLOS 2027 / ICLR 2026}
}

\author{%
  [Author Names Redacted for Review]\\[2pt]
  \small[Institution Redacted]\quad
  \small\texttt{[email redacted]}
}

\date{}

%% ════════════════════════════════════════════════════════════════════════════
\begin{document}
\maketitle
\thispagestyle{fancy}

%% ── Abstract ────────────────────────────────────────────────────────────────
\begin{abstract}
\noindent
Concurrency bugs in async and multi-threaded code cause severe production
failures, yet static detectors suffer high false-positive rates and LLMs lack
reliable reasoning about thread interactions.
Rule-based detectors such as NodeCB miss complex patterns; LLM-only approaches
(PCWMs) achieve only 75\% accuracy on race detection without explicit structural
analysis of the code.
ThreadLearn combines a 5-pattern static race detector with a Qwen2.5-Coder-1.5B
model fine-tuned via QLoRA on Hindsight Chain-of-Thought traces grounded in
static analysis outputs.
On [benchmark], ThreadLearn detects races with F1\,=\,X\% ($+$Y\% over
NodeCB baseline) and produces accepted fixes in Z\% of cases
($+$W\% over PCWMs baseline).
ThreadLearn is open-source at [repo]; the fine-tuning dataset and evaluation
harness are publicly available.
\end{abstract}

\noindent\textbf{Keywords:} concurrency bugs, race conditions, static analysis,
LLM fine-tuning, Chain-of-Thought, JavaScript

\bigskip\hrule\bigskip

%% ════════════════════════════════════════════════════════════════════════════
\section{Introduction}
\label{sec:intro}

\subsection{Problem Statement}

Asynchronous JavaScript (Node.js) powers the majority of modern web backends,
yet the language exposes concurrency primitives that are easy to misuse.
A landmark empirical study (NodeCB~\cite{nodecb2017}) examined 57 real-world
concurrency bugs across 53 open-source Node.js projects and found that
65\% are \emph{atomicity violations}, 30\% are \emph{order violations}, and
93\% cause severe failures including crashes, wrong database state, or hangs.

Despite this severity, the two dominant tool families remain inadequate in
complementary ways.
Rule-based detectors flag suspicious patterns but cannot suggest a fix and
accumulate false positives on idiomatic async code.
LLM-based approaches reason about code holistically but lack line-level causal
attribution and require large models (7B--32B parameters) to achieve competitive
accuracy.

\subsection{Gap Analysis}

We identify four concrete gaps in the state of the art:

\begin{description}[leftmargin=2em, labelwidth=2em]
  \item[\textbf{G1}] Rule-based detectors cannot fix bugs, only detect them.
  \item[\textbf{G2}] PCWMs~\cite{pcwms2026} achieves 75.2\% accuracy but does
    not exploit static structure and requires 27K training samples with a large
    base model.
  \item[\textbf{G3}] No existing system combines static analysis as a grounding
    signal for LLM reasoning.
  \item[\textbf{G4}] No end-to-end pipeline spans detection $\to$ explanation
    $\to$ fix for JavaScript (Node.js).
\end{description}

\subsection{Key Insight}

Static detector outputs---pattern identifier, line range, and causal
description---provide the missing \emph{grounding signal} for LLM concurrency
reasoning.
A small model (1.5B parameters) fine-tuned with grounded Hindsight
Chain-of-Thought (CoT) outperforms a large model prompted with generic CoT
because it has been explicitly trained to reason from structural evidence rather
than surface-level token patterns.

\subsection{Contributions}

This paper makes four contributions:

\begin{enumerate}[leftmargin=2em]
  \item \textbf{\texttt{race\_detector}}: a 5-pattern static analyzer for
    JavaScript (Node.js) integrated with an AST preprocessor (\S\ref{sec:design:detector}).
  \item \textbf{Hindsight CoT dataset}: a synthetic fine-tuning corpus of
    \texttt{(code, static\_output, CoT\_trace, fix)} tuples grounded on
    detector outputs (\S\ref{sec:design:finetuning}).
  \item \textbf{ThreadLearn pipeline}: an end-to-end detect $\to$ format $\to$
    RAG-augmented fix system exposed via a FastAPI server
    (\S\ref{sec:design:pipeline}).
  \item \textbf{Evaluation}: comparison against NodeCB and PCWMs baselines
    across detection F1, false-positive rate, fix acceptance rate, and latency
    (\S\ref{sec:eval}).
\end{enumerate}

\begin{figure}[ht]
  \centering
  \fbox{\parbox{0.85\linewidth}{\centering\vspace{4pt}
    \textbf{Figure 1 --- System Overview}\\[4pt]
    code input $\to$ AST preprocessor $\to$ race detector $\to$ report
    formatter $\to$ RAG pipeline $\to$ fine-tuned LLM $\to$ fix output\\[4pt]
    \textit{Baselines: NodeCB rule-only, PCWMs LLM-only (shown for comparison)}
  \vspace{4pt}}}
  \caption{ThreadLearn end-to-end system overview.}
  \label{fig:overview}
\end{figure}

%% ════════════════════════════════════════════════════════════════════════════
\section{Background and Motivation}
\label{sec:background}

\subsection{Concurrency Bug Taxonomy}

Table~\ref{tab:taxonomy} summarizes the concurrency bug types from the NodeCB
study and their JavaScript (Node.js) manifestations relevant to ThreadLearn.

\begin{table}[ht]
  \centering
  \caption{Concurrency bug taxonomy (NodeCB rates) with ThreadLearn examples.}
  \label{tab:taxonomy}
  \begin{tabular}{@{}llp{6cm}@{}}
    \toprule
    \textbf{Bug Type} & \textbf{Rate} & \textbf{Example} \\
    \midrule
    Atomicity violation & 65\% & Two interleaved DB calls cause lost update \\
    Order violation     & 30\% & Upload completes before create callback fires \\
    Starvation          & 5\%  & \texttt{setImmediate} starved by I/O queue \\
    \midrule
    \texttt{var i} + \texttt{setTimeout} & JS  & Closure captures wrong loop iteration value \\
    \bottomrule
  \end{tabular}
\end{table}

\subsection{Limitations of Prior Approaches}

\begin{table}[ht]
  \centering
  \caption{Prior approach limitations and key observations.}
  \label{tab:limits}
  \begin{tabular}{@{}lp{9.5cm}@{}}
    \toprule
    \textbf{Approach} & \textbf{Limitation / Observation} \\
    \midrule
    NodeCB study  & Regex-based detection; high false-positive rate; no fix generation \\
    PCWMs         & 75.2\% accuracy; does not exploit static structure;
                    requires large model (7B--32B) \\
    \midrule
    \textbf{O1}   & Static detector outputs (line range, pattern ID) provide
                    causal grounding missing from LLM reasoning \\
    \textbf{O2}   & A small model (1.5B) fine-tuned with grounded CoT
                    outperforms a large model with generic CoT \\
    \bottomrule
  \end{tabular}
\end{table}

\subsection{AST Preprocessing Foundation}

ThreadLearn's AST preprocessor exposes four functions used by both the race
detector and the RAG retrieval stage:

\begin{table}[ht]
  \centering
  \caption{AST preprocessor functions.}
  \label{tab:ast}
  \begin{tabular}{@{}llp{5cm}@{}}
    \toprule
    \textbf{Function} & \textbf{Purpose} & \textbf{Notes} \\
    \midrule
    \texttt{stripComments}       & Remove comments               & Regex-based; handles \texttt{//} and \texttt{/* */} \\
    \texttt{normalizeWhitespace} & Normalize whitespace          & Collapse multi-line blanks, trim trailing spaces \\
    \texttt{extractFunctions}    & Extract function/class list   & Arrow functions: known limitation \\
    \texttt{extract\_keywords}   & Top-20 identifier tokens      & Used for RAG retrieval \\
    \bottomrule
  \end{tabular}
\end{table}

%% ════════════════════════════════════════════════════════════════════════════
\section{Design}
\label{sec:design}

\begin{figure}[ht]
  \centering
  \fbox{\parbox{0.85\linewidth}{\centering\vspace{4pt}
    \textbf{Figure 2 --- Component Architecture}\\[4pt]
    AST preprocessor $\to$ race detector $\to$ report formatter
    $\to$ RAG pipeline $\to$ fine-tuned LLM
  \vspace{4pt}}}
  \caption{ThreadLearn five-component architecture with data flow.}
  \label{fig:arch}
\end{figure}

\subsection{System Architecture Overview}

ThreadLearn comprises five loosely coupled modules (Table~\ref{tab:components}):

\begin{table}[ht]
  \centering
  \caption{ThreadLearn system components and roles.}
  \label{tab:components}
  \begin{tabular}{@{}ll@{}}
    \toprule
    \textbf{Component} & \textbf{Role} \\
    \midrule
    \texttt{ast\_preprocessor} & Clean and normalize input code \\
    \texttt{race\_detector}    & Detect 5 JS race condition patterns \\
    \texttt{report\_formatter} & Map \texttt{pattern\_id} $\to$ severity + fix template \\
    \texttt{rag\_pipeline}     & Retrieve similar examples, augment LLM prompt \\
    Fine-tuned LLM             & Generate fix with CoT reasoning \\
    FastAPI server             & Expose \texttt{/analyze} endpoint \\
    \bottomrule
  \end{tabular}
\end{table}

\subsection{Static Race Detector}
\label{sec:design:detector}

The detector scans source files in $O(n)$ time and emits structured reports of
the form \texttt{\{pattern\_id, line\_range, description\}}.
Table~\ref{tab:patterns} lists all 5 patterns.
TypeScript files are treated as JavaScript.

\begin{table}[ht]
  \centering
  \caption{Race detection patterns in ThreadLearn (5 patterns, JavaScript only).}
  \label{tab:patterns}
  \begin{tabular}{@{}lll@{}}
    \toprule
    \textbf{Group} & \textbf{Pattern ID} & \textbf{Lang} \\
    \midrule
    \multirow{3}{*}{Closure/async}
      & \texttt{closure\_loop\_var}       & JS \\
      & \texttt{shared\_var\_settimeout}  & JS \\
      & \texttt{promise\_no\_await}       & JS \\
    \midrule
    \multirow{2}{*}{Shared state}
      & \texttt{concurrent\_write\_array} & JS \\
      & \texttt{counter\_no\_atomic}      & JS \\
    \bottomrule
  \end{tabular}
\end{table}

\paragraph{Example: closure loop var (JavaScript).}

\begin{lstlisting}[style=js, caption={Buggy: \texttt{var i} captured by reference.}]
for (var i = 0; i < 5; i++) {
  setTimeout(function() {
    console.log(i); // always prints 5
  }, 100);
}
\end{lstlisting}

\begin{lstlisting}[style=js, caption={Fixed: \texttt{let i} creates per-iteration binding.}]
for (let i = 0; i < 5; i++) {
  setTimeout(function() {
    console.log(i); // prints 0, 1, 2, 3, 4
  }, 100);
}
\end{lstlisting}


\subsection{Hindsight CoT Fine-Tuning}
\label{sec:design:finetuning}

Inspired by PCWMs~\cite{pcwms2026}, we construct training tuples:
\[
  (\textit{code},\; \textit{static\_output},\; \textit{CoT\_trace},\; \textit{fix})
\]
where \textit{static\_output} is the detector report and \textit{CoT\_trace}
is written by a teacher model \emph{after} the correct fix is known (``hindsight'').

\begin{table}[ht]
  \centering
  \caption{Fine-tuning configuration.}
  \label{tab:finetune}
  \begin{tabular}{@{}ll@{}}
    \toprule
    \textbf{Parameter} & \textbf{Value} \\
    \midrule
    Base model         & Qwen2.5-Coder-1.5B \\
    Method             & QLoRA: $r$=16, $\alpha$=32, 4-bit NF4 quantization \\
    Framework          & \texttt{SFTTrainer} (trl $\geq$ 1.5.0) \\
    Grounding signal   & \texttt{pattern\_id} + \texttt{line\_range} from detector \\
    \midrule
    \multicolumn{2}{@{}l@{}}{\textit{Rejected alternatives}} \\
    Full fine-tune     & Compute cost prohibitive \\
    RAG-only           & No explicit reasoning capability \\
    \bottomrule
  \end{tabular}
\end{table}

\subsection{RAG-Augmented Fix Pipeline}
\label{sec:design:pipeline}

\begin{enumerate}[leftmargin=2em]
  \item \textbf{Embed}: encode input code as a dense vector.
  \item \textbf{Retrieve}: find the $k$ most similar fixed examples from the RAG store.
  \item \textbf{Format}: \texttt{report\_formatter} maps \texttt{pattern\_id}
    $\to$ severity label and fix template string.
  \item \textbf{Generate}: fine-tuned LLM receives RAG context + detector
    report and produces a fix with CoT trace.
\end{enumerate}

\subsection{Rejected Design Alternatives}
\label{sec:design:rejected}

\begin{table}[ht]
  \centering
  \caption{Rejected design alternatives and justification.}
  \label{tab:rejected}
  \begin{tabular}{@{}lp{8cm}@{}}
    \toprule
    \textbf{Alternative} & \textbf{Reason Rejected} \\
    \midrule
    LLM-only detection     & High false-positive rate; no line-level attribution \\
    Static-only            & Detects bugs but cannot generate fixes \\
    Large model (7B+)      & Impractical for local deployment due to GPU memory \\
    AST-based JS detection & Requires Node.js runtime at analysis time \\
    \bottomrule
  \end{tabular}
\end{table}

%% ════════════════════════════════════════════════════════════════════════════
\section{Implementation}
\label{sec:impl}

\begin{table}[ht]
  \centering
  \caption{Implementation summary.}
  \label{tab:impl}
  \begin{tabular}{@{}ll@{}}
    \toprule
    \textbf{Item} & \textbf{Detail} \\
    \midrule
    Language / framework     & Python 3.10+, FastAPI \\
    Total LOC                & $\approx$1,200 \\
    \texttt{race\_detector.py} & 320 LOC, 10 detectors, $O(n)$ scan \\
    Model adapter            & LoRA SafeTensors at
                               \texttt{./threadlearn-qwen2.5-coder-1.5b/final\_adapter} \\
    Infrastructure           & Docker: Redis cache (localhost:6379), Atlas MongoDB \\
    Concurrency control      & Semaphore in FastAPI server \\
    Key engineering note     & Graceful fallback if \texttt{ast\_preprocessor}
                               import fails (CI isolation) \\
    \bottomrule
  \end{tabular}
\end{table}

%% ════════════════════════════════════════════════════════════════════════════
\section{Evaluation}
\label{sec:eval}

\subsection{Experimental Setup}

\begin{table}[ht]
  \centering
  \caption{Experimental setup.}
  \label{tab:setup}
  \begin{tabular}{@{}ll@{}}
    \toprule
    \textbf{Item} & \textbf{Detail} \\
    \midrule
    Baseline 1     & NodeCB-style rule-only detector \\
    Baseline 2     & PCWMs-style LLM-only (Qwen2.5-Coder-1.5B, no fine-tune) \\
    Our system     & ThreadLearn (static detector + fine-tuned LLM) \\
    Datasets       & 20 synthetic test cases + [real-world dataset TBD] \\
    Metrics        & Detection F1, FPR, Fix Acceptance Rate, Latency (ms/req) \\
    \bottomrule
  \end{tabular}
\end{table}

\subsection{End-to-End Detection Comparison}
\label{sec:eval:detection}

\begin{table}[ht]
  \centering
  \caption{Detection results across three methods (values TBD pending AI1-07 fine-tuning run).}
  \label{tab:detection}
  \begin{tabular}{@{}lcccc@{}}
    \toprule
    \textbf{Method} & \textbf{Prec.} & \textbf{Recall} & \textbf{F1} & \textbf{FPR} \\
    \midrule
    NodeCB rule-only                & -- & -- & -- & -- \\
    PCWMs LLM-only                  & -- & -- & -- & -- \\
    \textbf{ThreadLearn (ours)}     & -- & -- & -- & -- \\
    \midrule
    \multicolumn{5}{@{}l@{}}{\textit{JS subset}} \\
    \quad ThreadLearn               & -- & -- & -- & -- \\
    \bottomrule
  \end{tabular}
\end{table}

\begin{figure}[ht]
  \centering
  \fbox{\parbox{0.85\linewidth}{\centering\vspace{4pt}
    \textbf{Figure 3 --- F1 Bar Chart}\\[4pt]
    3 methods, JavaScript (Node.js) benchmark\\[2pt]
    \textit{(placeholder --- data TBD)}
  \vspace{4pt}}}
  \caption{F1 comparison: ThreadLearn vs.\ baselines.}
  \label{fig:f1}
\end{figure}

\subsection{Fix Quality Evaluation}
\label{sec:eval:fix}

\begin{table}[ht]
  \centering
  \caption{Fix acceptance rate and case studies.}
  \label{tab:fix}
  \begin{tabular}{@{}lp{9cm}@{}}
    \toprule
    \textbf{Metric / Case} & \textbf{Detail} \\
    \midrule
    Fix acceptance rate & Human evaluation or compile + test pass \\
    Case 1 (JS)         & \texttt{var i} closure $\to$ \texttt{let i} \\
    PCWMs comparison    & Original PCWMs: +2.7\%--11.1\% with world model feedback \\
    \bottomrule
  \end{tabular}
\end{table}

\begin{figure}[ht]
  \centering
  \fbox{\parbox{0.85\linewidth}{\centering\vspace{4pt}
    \textbf{Figure 4 --- Fix Acceptance Rate}\\[4pt]
    ThreadLearn vs.\ LLM-only vs.\ rule-only (no fix)\\[2pt]
    \textit{(placeholder --- data TBD)}
  \vspace{4pt}}}
  \caption{Fix acceptance rate comparison.}
  \label{fig:fix}
\end{figure}

\subsection{Ablation Study}
\label{sec:eval:ablation}

\begin{table}[ht]
  \centering
  \caption{Ablation: per-component contribution (values TBD pending fine-tuning run).}
  \label{tab:ablation}
  \begin{tabular}{@{}lp{8cm}@{}}
    \toprule
    \textbf{Component Removed} & \textbf{Expected Impact} \\
    \midrule
    Static grounding    & F1 drops X\% \\
    RAG retrieval       & Fix quality drops Y\% \\
    Hindsight CoT       & Accuracy drops Z\% (validates PCWMs finding) \\
    QLoRA $r$=8 vs.\ 16 & Accuracy vs.\ GPU memory tradeoff \\
    \bottomrule
  \end{tabular}
\end{table}

\subsection{Scalability and Latency}
\label{sec:eval:latency}

\begin{table}[ht]
  \centering
  \caption{Latency and throughput from Locust load test (AI2-10 incomplete).}
  \label{tab:latency}
  \begin{tabular}{@{}ll@{}}
    \toprule
    \textbf{Metric} & \textbf{Value} \\
    \midrule
    Requests/second  & [TBD] \\
    P95 latency      & [TBD] \\
    Semaphore impact & [TBD --- effect on throughput] \\
    \bottomrule
  \end{tabular}
\end{table}

%% ════════════════════════════════════════════════════════════════════════════
\section{Related Work}
\label{sec:related}

\begin{table}[ht]
  \centering
  \caption{Related work: how ThreadLearn differs from prior approaches.}
  \label{tab:related}
  \begin{tabular}{@{}p{3.2cm}p{4.5cm}p{4.5cm}@{}}
    \toprule
    \textbf{Group} & \textbf{Representative Works} & \textbf{How ThreadLearn Differs} \\
    \midrule
    Static race detectors
      & ThreadSanitizer, Helgrind, ESLint async rules
      & Detect without fixing; no LLM reasoning \\
    \addlinespace
    Empirical bug studies
      & NodeCB~\cite{nodecb2017}, Android concurrency bugs
      & Taxonomy source; no automated fix \\
    \addlinespace
    LLM for code analysis
      & PCWMs~\cite{pcwms2026}, CodeBERT, CodeT5
      & LLM-only; does not exploit static structure \\
    \addlinespace
    Fine-tuning for code
      & QLoRA~\cite{qlora}, SFT on code, PEFT
      & Training methodology; not targeting race detection \\
    \bottomrule
  \end{tabular}
\end{table}

\paragraph{Method comparison matrix.}

\begin{table}[ht]
  \centering
  \caption{Method $\times$ feature comparison matrix.}
  \label{tab:matrix}
  \begin{tabular}{@{}lcccc@{}}
    \toprule
    \textbf{Method} & \textbf{Detect} & \textbf{Fix} & \textbf{JS} & \textbf{Fine-tuned} \\
    \midrule
    ThreadSanitizer      & $\checkmark$ & $\times$ & $\times$     & $\times$ \\
    ESLint async         & $\checkmark$ & $\times$ & $\checkmark$ & $\times$ \\
    NodeCB (rules)       & $\checkmark$ & $\times$ & $\checkmark$ & $\times$ \\
    PCWMs (LLM-only)     & $\checkmark$ & $\checkmark$ & $\checkmark$ & $\checkmark$ \\
    \textbf{ThreadLearn} & $\checkmark$ & $\checkmark$ & $\checkmark$ & $\checkmark$ \\
    \bottomrule
  \end{tabular}
\end{table}

%% ════════════════════════════════════════════════════════════════════════════
\section{Conclusion}
\label{sec:conclusion}

Concurrency bugs remain a leading cause of production failures in async and
multi-threaded applications, yet no existing tool both detects and fixes them
reliably.
ThreadLearn combines a 5-pattern static race detector with a Hindsight CoT
fine-tuned Qwen2.5-Coder-1.5B model to provide an end-to-end detection and
fix generation pipeline for JavaScript (Node.js).
ThreadLearn achieves F1\,=\,X\% on race detection ($+$Y\% over NodeCB) and fix
acceptance rate Z\% ($+$W\% over PCWMs baseline), demonstrating that static
grounding significantly improves LLM concurrency reasoning.

\paragraph{Future work.}
Expand to TypeScript AST (replace regex fallback);
add MPI/CUDA patterns (PCWMs direction).

%% ════════════════════════════════════════════════════════════════════════════
\begin{thebibliography}{9}

\bibitem{nodecb2017}
  T.\ Xu, X.\ Jin, P.\ Huang, Y.\ Zhou, S.\ Lu, L.\ Wang, and
  K.\ V.\ Vishwanath.
  \newblock {NodeCB}: Understanding Concurrency Bugs in Node.js.
  \newblock In \emph{Proc.\ 32nd IEEE/ACM ASE}, pp.\ 87--98, 2017.

\bibitem{pcwms2026}
  [Author(s) redacted].
  \newblock {PCWMs}: Prompting with Chain-of-Thought World Models for
  Concurrency Bug Detection.
  \newblock \emph{arXiv:2604.20926v2}, 2026.

\bibitem{qlora}
  T.\ Dettmers, A.\ Pagnoni, A.\ Holtzman, and L.\ Zettlemoyer.
  \newblock {QLoRA}: Efficient Finetuning of Quantized LLMs.
  \newblock In \emph{Proc.\ NeurIPS}, 2023.

\bibitem{threadsanitizer}
  K.\ Serebryany and T.\ Iskhodzhanov.
  \newblock {ThreadSanitizer}: Data Race Detection in Practice.
  \newblock In \emph{Proc.\ WBIA}, pp.\ 62--71, 2009.

\bibitem{qwen25coder}
  Qwen Team.
  \newblock {Qwen2.5-Coder} Technical Report.
  \newblock \emph{arXiv:2409.12186}, 2024.

\end{thebibliography}

\end{document}
